% GENERATED FILE. Edit canonical lessons, then run npm run generate:books.
% canonical-source-hash: fnv1a64:0083c19e2718a6a1
% canonical-lessons: ZH-C85-zai, ZH-W85-zai-at, ZH-C85-nar, ZH-C85-libian, ZH-C85-waibian, ZH-C85-practice

\chapter{Where Things Are}
\label{ch:zh-where-things-are}
\hlchaptermodality{85}

\begin{chapteropening}
\textbf{By the end of this chapter:} \emph{I can say where people and things are, and write \zh{在}.}

Saying where people and things are with \zh{在}, here and there, inside and outside, built on one new written character: \zh{在}.
\end{chapteropening}

\section[\texorpdfstring{\zh{在} (zài)}{在 (zài)}]{\zh{在} (zài) — at, in; to be somewhere}

\label{lesson:ZH-C85-zai}



\begin{quote}
A new chapter, and it opens with the word that puts people and things in places.

\begin{itemize}
\raggedright
  \item \emph{Recall:} say \emph{zhǔnbèi}, to get ready — \textbf{R1}, the last word you met
  \item \emph{Recall:} say \zh{这儿} — \emph{zhèr}, here — \textbf{R3}, durable, from the chapter on home
  \item \emph{Recall:} say \emph{pángbiān}, beside — \textbf{R2}, a few chapters back
\end{itemize}
\end{quote}

\subsection*{You'll want to know: \zh{在}}
\textbf{\zh{在}} — \emph{zài} — \textquotedblleft{}to be at; in, at\textquotedblright{}.

\begin{quote}
\textbf{\zh{妈妈在家吗}} \emph{māma zài jiā ma} — Is Mum at home?
\end{quote}

\begin{quote}
\textbf{\zh{我在这儿}} \emph{wǒ zài zhèr} — I'm here.
\end{quote}

The next lesson puts the character in your hand.

\begin{cousinweb}
\textbf{\zh{在}} (\emph{zài}) is inherited from Middle Chinese \emph{dzojX} (Guangyun, in Baxter's transcription), glossed 'to dwell' and 'to exist'.
\end{cousinweb}

\begin{grammarlens}[title={What you've built}]
The place comes straight after \zh{在}: \zh{妈妈}, \zh{在}, \zh{家}. English needs \emph{is} and \emph{at}; Mandarin needs \zh{在} alone.

That is the whole sentence pattern of the chapter: a person or a thing, then \zh{在}, then the place.
\end{grammarlens}

\subsection*{Your turn}
\begin{itemize}
\raggedright
  \item \emph{Say it:} \emph{zài}
  \item \emph{Say it:} \zh{妈妈在家吗} — \emph{māma zài jiā ma} — Is Mum at home?
  \item \emph{Say it:} \zh{我在这儿} — \emph{wǒ zài zhèr} — I'm here.
  \item \emph{Recall:} say \emph{zhǔnbèi} and \emph{pángbiān}, then \emph{zài} again
\end{itemize}

\begin{tcolorbox}[breakable,colback=teal!4,colframe=teal!35!black,title={Before you move on}]
What does \zh{在} mean? (\textquotedblleft{}To be at; in, at\textquotedblright{}.) Where does the place go? (\textbf{Straight after \zh{在}}.)
\end{tcolorbox}
\section[\texorpdfstring{\zh{在} (zài)}{在 (zài)}]{\zh{在} — the top of \zh{有}, a vertical, then \zh{土}}

\label{lesson:ZH-W85-zai-at}



\begin{quote}
The word first, then the character inside it.

\begin{itemize}
\raggedright
  \item \emph{Recall:} say \emph{zài} — \textbf{R1}, one lesson back, before it decays
\end{itemize}

\begin{quote}
\zh{有} \zh{里} \zh{十}
\end{quote}

\emph{Read it:} each one aloud, with nothing said first — \emph{yǒu}, \emph{lǐ}, \emph{shí}
\end{quote}

\subsection*{Script — the top of \zh{有}, a vertical, then \zh{土}}
\hlblockfigure{\detokenize{figures/ZH-W85-zai-at-filmstrip.pdf}}{How \zh{在} is written, stroke by stroke}

\begin{quote}
\zh{在}
\end{quote}

\textbf{Three pieces, and two of them you already write.}

\textbf{The top of \zh{有}, two strokes}: a long horizontal, then a long left-falling stroke from above it, down to the lower left.

\textbf{A short vertical} under them, on the left, down from the falling stroke.

\textbf{\zh{土}, three strokes}: a short horizontal, a vertical down through it, and a long bar along the foot. It is \zh{十} standing on a bar.

2 + 1 + 3 = six.

\subsection*{Your turn}
\begin{itemize}
\raggedright
  \item \emph{Write it:} \zh{在}, all six, with the model in front of you
  \item \emph{Cover:} write it again from memory, and count the pieces, not the strokes
\end{itemize}

\begin{tcolorbox}[breakable,colback=teal!4,colframe=teal!35!black,title={Before you move on}]
How many strokes in \zh{在}? (\textbf{Six}.) Which character you know opens the same way? (\textbf{\zh{有}}.) Which bar is longest? (\textbf{The last, along the foot}.)

Source: \href{https://raw.githubusercontent.com/chanind/hanzi-writer-data/68d10a4b21150cae5e1ebbd223eed289cf32d90c/data/\%E5\%9C\%A8.json}{Hanzi Writer Data \zh{在}, pinned stroke paths and medians}.
\end{tcolorbox}
\section[\texorpdfstring{\zh{那儿} (nàr)}{那儿 (nàr)}]{\zh{那儿} (nàr) — there}

\label{lesson:ZH-C85-nar}



\begin{quote}
Before the new one: the character from the last lesson, then a word it pairs with.

\begin{itemize}
\raggedright
  \item \emph{Read it:} \zh{在家}, aloud — \emph{zài jiā}, at home
  \item \emph{Recall:} say \zh{这儿} — \emph{zhèr}, here
\end{itemize}
\end{quote}

\subsection*{You'll want to know: \zh{那儿}}
\textbf{\zh{那儿}} — \emph{nàr} — \textquotedblleft{}there\textquotedblright{}.

\begin{quote}
\textbf{\zh{他在那儿}} \emph{tā zài nàr} — He's over there.
\end{quote}

\begin{quote}
\textbf{\zh{我在这儿}} \emph{wǒ zài zhèr} — I'm here.
\end{quote}

\begin{cousinweb}
\textbf{\zh{那儿}} (\emph{nàr}) is built from \zh{那} 'that' and \zh{儿}, the ending that runs into the syllable before it as an r: the same two pieces as \zh{这儿}. The Guangyun records \zh{那} as \emph{na} (Guangyun, in Baxter's transcription), glossed 'how, what'; how it came to mean 'that' is not settled.
\end{cousinweb}

\begin{grammarlens}[title={What you've built}]
\zh{这儿} and \zh{那儿}, here and there, answer most of the place questions this chapter asks. Both take the r ending, and both follow \zh{在}.
\end{grammarlens}

\subsection*{Your turn}
Say these aloud:

\begin{itemize}
\raggedright
  \item \emph{nàr}
  \item \emph{nàr}, once more
  \item \zh{他在那儿} — \emph{tā zài nàr} — He's over there.
  \item \zh{我在这儿} \zh{你在那儿} — \emph{wǒ zài zhèr, nǐ zài nàr} — I'm here, you're there.
\end{itemize}

\begin{tcolorbox}[breakable,colback=teal!4,colframe=teal!35!black,title={Before you move on}]
What does \zh{那儿} mean? (\textquotedblleft{}There\textquotedblright{}.) And here? (\textbf{\zh{这儿}}.) Say both once more.
\end{tcolorbox}
\section[lǐbian]{lǐbian — inside}

\label{lesson:ZH-C85-libian}



\begin{quote}
Before the new one: last lesson's word, and a character you write.

\begin{itemize}
\raggedright
  \item \emph{Recall:} say \zh{那儿} — \emph{nàr}, there
  \item \emph{Read it:} \zh{里}, aloud — \emph{lǐ}, inside
  \item \emph{Recall:} say \emph{chéngshì}, a city
\end{itemize}
\end{quote}

\subsection*{You'll want to know: lǐbian}
\emph{lǐbian} — \textquotedblleft{}inside; in\textquotedblright{}.

\begin{quote}
\emph{shū zài bāo lǐbian} — The book is in the bag.
\end{quote}

Its second character comes later; for now it is a word for the ear.

\begin{cousinweb}
\emph{lǐbian} is built from lǐ 'inside', the \zh{里} you write, and bian 'side', the syllable that ends \emph{zuǒbian} and \emph{pángbiān}. Both are inherited from Middle Chinese, \emph{liX} and \emph{pen} (Guangyun, in Baxter's transcription), glossed 'the inside of a garment' and 'an edge; a side'.
\end{cousinweb}

\begin{grammarlens}[title={What you've built}]
\emph{lǐbian} comes after the thing whose inside it names: \emph{bāo lǐbian}, inside the bag. Put \zh{在} in front and it is a place: \emph{zài bāo lǐbian}.

\zh{里} on its own does the same job in a quicker sentence, as in \emph{shān lǐ}. \emph{lǐbian} can also stand alone: \emph{zài lǐbian}, inside.
\end{grammarlens}

\subsection*{Your turn}
\begin{itemize}
\raggedright
  \item \emph{Say it:} \emph{lǐbian}
  \item \emph{Say it:} \emph{lǐbian}, once more
  \item \emph{Say it:} \emph{shū zài bāo lǐbian} — The book is in the bag.
  \item \emph{Say it:} \emph{tā zài lǐbian} — He's inside.
  \item \emph{Recall:} say \emph{nàr} and \emph{chéngshì}, then \emph{lǐbian} again
\end{itemize}

\begin{tcolorbox}[breakable,colback=teal!4,colframe=teal!35!black,title={Before you move on}]
What does \emph{lǐbian} mean? (\textquotedblleft{}Inside; in\textquotedblright{}.) Which written character does it open with? (\textbf{\zh{里}}.)
\end{tcolorbox}
\section[wàibian]{wàibian — outside}

\label{lesson:ZH-C85-waibian}



\begin{quote}
Before the new one: say the word for inside, then two words you have had for a while.

\begin{itemize}
\raggedright
  \item \emph{Recall:} say \emph{lǐbian}, inside — \textbf{R1}, one lesson back
  \item \emph{Recall:} say \emph{wàiguó}, abroad — \textbf{R2}, a few chapters back
  \item \emph{Recall:} say \emph{zuǒbian} and \emph{yòubian}, left and right — \textbf{R2}
\end{itemize}
\end{quote}

\subsection*{You'll want to know: wàibian}
\emph{wàibian} — \textquotedblleft{}outside\textquotedblright{}.

\begin{quote}
\emph{tā zài wàibian} — She's outside.
\end{quote}

Its characters come later; for now it is a word for the ear.

\begin{cousinweb}
\emph{wàibian} is built from wài 'outside', the wài of \emph{wàiguó}, and bian 'side'. Both are inherited from Middle Chinese, \emph{ngwajH} and \emph{pen} (Guangyun, in Baxter's transcription), \emph{ngwajH} glossed 'the outside' and 'far'.
\end{cousinweb}

\begin{grammarlens}[title={What you've built}]
\emph{lǐbian} and \emph{wàibian}, inside and outside, beside the pair you already own, \emph{zuǒbian} and \emph{yòubian}: one ending, bian, makes all four, and all four follow \zh{在}.
\end{grammarlens}

\subsection*{Your turn}
\begin{itemize}
\raggedright
  \item \emph{Say it:} \emph{wàibian}
  \item \emph{Say it:} \emph{wàibian}, once more
  \item \emph{Say it:} \emph{tā zài wàibian} — She's outside.
  \item \emph{Say it:} \emph{māo zài wàibian} — The cat's outside.
  \item \emph{Recall:} say \emph{lǐbian} and \emph{wàiguó}, then \emph{wàibian} again
\end{itemize}

\begin{tcolorbox}[breakable,colback=teal!4,colframe=teal!35!black,title={Before you move on}]
What does \emph{wàibian} mean? (\textquotedblleft{}Outside\textquotedblright{}.) And inside? (\emph{\textbf{lǐbian}}.)
\end{tcolorbox}
\section[Practice]{\zh{我在这儿} — \zh{你在那儿}}

\label{lesson:ZH-C85-practice}



\begin{quote}
The chapter closes here, with two people placing themselves and their things.

\begin{itemize}
\raggedright
  \item \emph{Recall:} say \emph{wàibian}, outside — \textbf{R1}, one lesson back
  \item \emph{Recall:} say \emph{yìqǐ}, together — \textbf{R2}, the first real retrieval
  \item \emph{Recall:} say \emph{jiànmiàn}, to meet up — \textbf{R2}
\end{itemize}
\end{quote}

\subsection*{The exchange}
\begin{quote}
\textbf{A:} \zh{我在家} \zh{你呢}
\end{quote}

\begin{quote}
\emph{wǒ zài jiā, nǐ ne} — I'm at home. And you?
\end{quote}

\begin{quote}
\textbf{B:} \emph{wǒ zài xuéxiào} — I'm at school.
\end{quote}

\begin{quote}
\textbf{A:} \zh{你的书呢}
\end{quote}

\begin{quote}
\emph{nǐ de shū ne} — Where's your book?
\end{quote}

\begin{quote}
\textbf{B:} \emph{zài bāo lǐbian} — In my bag.
\end{quote}

\begin{quote}
\textbf{A:} \emph{māo ne} — And the cat?
\end{quote}

\begin{quote}
\textbf{B:} \emph{māo zài wàibian} — The cat's outside.
\end{quote}

\begin{quote}
\textbf{A:} \zh{他呢}
\end{quote}

\begin{quote}
\emph{tā ne} — And him?
\end{quote}

\begin{quote}
\textbf{B:} \zh{他在那儿}
\end{quote}

\begin{quote}
\emph{tā zài nàr} — He's over there.
\end{quote}

\emph{nǐ de shū ne} asks where the book is: \zh{呢} after a thing, with nothing else, means \emph{and where is it?} The answer drops the book and starts at \zh{在}.

\begin{tcolorbox}[breakable,skin=enhanced,colback=orange!6,colframe=orange!45!black,boxrule=0.5pt,arc=1mm,left=6pt,right=6pt,top=4pt,bottom=4pt,fonttitle=\bfseries,title={Read this}]
\emph{Read it:} the exchange through once without stopping, then answer the question below
\end{tcolorbox}

\begin{grammarlens}[title={What you've built — one pattern, four places}]
One pattern, a person or a thing, then \zh{在}, then the place:

\begin{itemize}
\raggedright
  \item \textbf{\zh{我在这儿}} \emph{wǒ zài zhèr} — I'm here.
  \item \textbf{\zh{他在那儿}} \emph{tā zài nàr} — he's there.
  \item \emph{shū zài bāo lǐbian} — the book is in the bag.
  \item \emph{māo zài wàibian} — the cat is outside.
\end{itemize}

\textbf{Listen.} \emph{Nǐ de shǒujī ne?}

\textbf{Speak.} Place yourself, then somebody you know, then your phone.

\emph{Write it:} \zh{在} — the top of \zh{有}, a short vertical, then \zh{十} on a long bar
\end{grammarlens}

\subsection*{Your turn}
\begin{itemize}
\raggedright
  \item \emph{Say it:} \zh{我在家} \zh{你呢} — \emph{wǒ zài jiā, nǐ ne} — I'm at home. And you?
  \item \emph{Say it:} \zh{你的书呢} — \emph{nǐ de shū ne} — Where's your book?
  \item \emph{Say it:} the whole exchange, both parts
  \item \emph{Write it:} \zh{我在这儿}
\end{itemize}

\begin{tcolorbox}[breakable,colback=teal!4,colframe=teal!35!black,title={Before you move on}]
Say you are at home. (\textbf{\zh{我在家}}.) Say the book is in the bag. (\emph{shū zài bāo lǐbian}.) Ask where somebody's book is, with three characters. (\textbf{\zh{你的书呢}}.)
\end{tcolorbox}
